\documentclass[11pt,a4paper]{article}
\usepackage[utf8]{inputenc}
\usepackage[french]{babel}
\usepackage[T1]{fontenc}
\usepackage{amsmath,amsfonts,amssymb}
\usepackage{geometry}
\geometry{margin=1.5cm}
\usepackage{tcolorbox}
\usepackage{graphicx}

\title{\textbf{Entraînement Brevet : Géométrie}}
\author{Mathématiques --- Thales \& Pythagore}
\maDate{}

\begin{document}

    \maketitle
    
    \begin{tcolorbox}[colback=blue!5!white,colframe=blue!75!black,title=Rappels de cours]
    \begin{itemize}
        \item \textbf{Le théorème de Pythagore :} Sert à calculer une longueur dans un triangle \textbf{rectangle} ($a^2 + b^2 = c^2$). La réciproque sert à prouver qu'un triangle \textbf{est} rectangle.
        \item \textbf{Le théorème de Thalès :} Sert à calculer des longueurs dans une configuration de droites \textbf{parallèles}. La réciproque sert à prouver que des droites \textbf{sont} parallèles.
    \end{itemize}
    \end{tcolorbox}
    
    \section*{Exercice 1 : Autour du triangle rectangle (Pythagore)}
    
    Soit $ABC$ un triangle tel que $AB = 4,5$ cm, $AC = 6$ cm et $BC = 7,5$ cm.
    On considère également un point $D$ tel que le triangle $ABD$ soit rectangle en $A$ avec $AD = 2,4$ cm.
    
    \begin{enumerate}
        \item Démontrer que le triangle $ABC$ est un triangle rectangle en $A$.
        \item Calculer la longueur du segment $[BD]$. Donner la valeur exacte.
        \item \textit{Question bonus :} Calculer l'aire du triangle $ABC$.
    \end{enumerate}
    
    \vspace{1cm}
    
    
    \section*{Exercice 2 : La tyrolienne (Thalès)}
    
        Un moniteur de parc d'attraction installe une tyrolienne entre deux poteaux verticaux $[AB]$ et $[CD]$ fixés au sol. Le sol est représenté par la droite $(BD)$.
        On donne les informations suivantes :
        \begin{itemize}
            \item Les poteaux sont perpendiculaires au sol : $(AB) \perp (BD)$ et $(CD) \perp (BD)$.
            \item $AB = 15$ m et $CD = 6$ m.
        \end{itemize}
        
        On considère la figure simplifiée ci-dessous (qui n'est pas à l'échelle) où les droites $(AC)$ et $(BD)$ se coupent en $E$. On donne $EB = 20$ m, $ED = 8$ m et $EA = 25$ m.
        
        \begin{enumerate}
            \item Montrer que les droites $(AB)$ et $(CD)$ sont parallèles.
            \item Calculer la longueur $EC$.
            \item Un troisième poteau $[FG]$ de $7$ m de haut est placé de telle sorte que $G$ soit sur $[ED]$ avec $EG = 5,6$~m. Les droites $(FG)$ et $(AB)$ sont-elles parallèles ? Justifier.
        \end{enumerate}
    
    \begin{center}
    \begin{tikzpicture}[scale=0.3] % Échelle pour que la figure tienne sur la largeur
        
        % Le sol (droite horizontale)
        \draw[thick] (-2,0) -- (15,0);
        \node[below] at (6,0) {Sol};
    
        % Les poteaux (segments verticaux)
        \draw[ultra thick, brown] (0,0) -- (0,15) node[midway, left, black] {15 m}; % Poteau AB
        \draw[ultra thick, brown] (16,0) -- (16,8.3) node[midway, right, black] {6 m}; % Poteau CD
    
        
        % Prolongement pour la configuration de Thalès (point E)
        % On calcule le point d'intersection E des droites (AC) et (BD)
        % Avec AB=15 et CD=6, le point E est à droite (configuration papillon décalée)
        % Ou plus simplement, on trace la configuration classique demandée dans l'énoncé :
        \coordinate (E) at (36,0); % Point E décalé pour la perspective de calcul
        \draw[thick] (0,0) -- (36,0);
        \draw (0,15) -- (36,0);
    
        % Nom des points
        \node[below] at (0,0) {$B$};
        \node[above] at (0,15) {$A$};
        \node[below] at (16,0) {$D$};
        \node[above right] at (16,8) {$C$}; % Ajusté selon la pente AC
        \node[below right] at (36,0) {$E$};
    
        % Symboles d'angle droit
        \draw (0,0.5) -- (0.5,0.5) -- (0.5,0);
        \draw (16,0.5) -- (16.5,0.5) -- (16.5,0);
    
        % Troisième poteau (Question 3)
        \draw[ultra thick, green!50!black] (22,0) -- (21,6.2); % Poteau FG indicatif
        \node[below] at (22,0) {$G$};
        \node[above] at (21.6,6.6) {$F$};
        
    \end{tikzpicture}
    \end{center}


\newpage


    \begin{tcolorbox}[colback=yellow!5!white,colframe=yellow!75!black,title=Aide Mémoire]
    \begin{itemize}
        \item \textbf{Triangles semblables :} Deux triangles sont semblables si leurs côtés sont proportionnels.
        \item \textbf{Vitesse :} $v = \frac{d}{t}$. Attention à bien convertir les unités (mètres en kilomètres et minutes en heures).
    \end{itemize}
    \end{tcolorbox}
    
       
    \section*{Exercice 3 : Le défi du VTT (Triangles semblables et Vitesse)}
    
    Un cycliste s'entraîne sur un parcours représenté par la figure ci-dessous. 
    Les points $A$, $S$ et $C$ sont alignés, ainsi que les points $B$, $S$ et $D$.
    On donne les dimensions suivantes : 
    $AS = 150$ m, $BS = 200$ m, $AB = 250$ m.
    $SC = 450$ m et $SD = 600$ m.

    
    \begin{center}
    \begin{tikzpicture}[scale=0.015] % Échelle adaptée pour que 600 unités tiennent sur la page
    
        % Coordonnées des points
        \coordinate (S) at (0,0);
        \coordinate (A) at (-150, 80);
        \coordinate (B) at (-200, -120);
        
        % Points alignés (Rapport 3 pour SCD)
        \coordinate (C) at (450, -240); % S est le milieu ? Non, alignés : C est de l'autre côté de S
        \coordinate (D) at (600, 360);
    
        % Tracé des triangles
        \draw[thick, fill=blue!5] (S) -- (A) -- (B) -- cycle;
        \draw[thick, fill=red!5] (S) -- (C) -- (D) -- cycle;
    
        % Marquage des points
        \node[above] at (A) {$A$};
        \node[below] at (B) {$B$};
        \node[below left] at (S) {$S$};
        \node[below] at (C) {$C$};
        \node[above] at (D) {$D$};
    
        % Affichage des longueurs (optionnel, on peut les laisser dans l'énoncé)
        \node[above left, sloped] at (-75, 40) {\small 150 m};
        \node[below left, sloped] at (-100, -60) {\small 200 m};
        \node[above, sloped] at (-175, -20) {\small 250 m};
        
        \node[below right, sloped] at (225, -120) {\small 450 m};
        \node[above right, sloped] at (300, 180) {\small 600 m};
        \node[right] at (525, 60) {?}; % Pour CD
    
        % Indication des angles opposés par le sommet
        \draw[red, thick] (-15, 8) arc (150:230:15);
        \draw[red, thick] (15, -8) arc (-30:50:15);
    
    \end{tikzpicture}
    \end{center}
    
    
    
    \begin{enumerate}
        \item Démontrer que les triangles $ABS$ et $SCD$ sont semblables.
        \item En déduire le rapport d'agrandissement pour passer du triangle $ABS$ au triangle $SCD$.
        \item Calculer la longueur $CD$ du deuxième tronçon du parcours.
        \item Le cycliste parcourt la distance totale $AB + BS + SD + DC$ en exactement 6 minutes.
        \begin{enumerate}
            \item Montrer que la distance totale parcourue est de $1\,850$ mètres.
            \item Calculer sa vitesse moyenne en m/min.
            \item Convertir cette vitesse en km/h. Est-il au-dessus de la limite de $20$ km/h autorisée sur ce sentier ?
        \end{enumerate}
    \end{enumerate}
    


\end{document}